\section{Serverless Min=0 Architecture \& Preemptible Spot GPU/TPU Economics}

Beyond physical fidelity, the transition of Earth System Digital Twins from academic benchmarks to production-grade operational deployments is heavily constrained by compute economics. Traditional atmospheric models scale via tightly-coupled MPI (Message Passing Interface) over dedicated HPC clusters. Reserving a traditional 128-core AMD EPYC / A100 node typically incurs static costs of $\$15$--$\$35$ per hour, operating 24/7 with substantial idle burn between forecast cycles.

LeanFlow radically alters this paradigm by decoupling the dual-scale execution into stateless, verifiable functional blocks that natively target serverless container orchestration architectures (e.g., Google Cloud Run Jobs, Vertex AI Serverless, Modal). This architecture allows the computational fabric to dynamically scale from zero to thousands of concurrent preemptible spot workers and back to zero instantly upon completion of the time-step block.

\subsection{Cost Profiling and Spot Economics}

Our empirical 10-minute benchmark executed the $128 \times 256$ grid solver utilizing preemptible Google Cloud spot instances:
\begin{itemize}
    \item \textbf{Spot Cloud TPU v5e (16GB HBM):} $\$0.36$/hour.
    \item \textbf{Spot Nvidia L4 GPU (24GB VRAM):} $\$0.20$/hour.
\end{itemize}

Figure~\ref{fig:serverless_economics}a demonstrates the cumulative execution cost over the 600-second run. While a traditional dedicated HPC cluster node accrued a static continuous cost leading to a high execution overhead, the serverless Spot L4 and TPU instances achieved up to a $98.9\%$ reduction in raw compute costs. Crucially, as highlighted in the scaling analysis, the instant scale-to-zero property eliminates idle burn rates completely ($0.00$/hr) when unallocated.

\begin{figure}[H]
    \centering
    \includegraphics[width=\textwidth]{figures/figure_serverless_spot_scaling.pdf}
    \caption{Serverless spot economics and spatial grid scaling. (a) Cumulative 600-second execution cost comparing traditional dedicated HPC nodes with Cloud On-Demand L4 and Serverless Spot TPU/GPU models. Note the instant scale-to-zero property upon completion. (b) Computational throughput (steps per second) scaling across grid resolutions ($32^2$ to $512^2$), highlighting the Neural-FGMRES accelerator performance against a standard 128-core CPU WRF baseline.}
    \label{fig:serverless_economics}
\end{figure}

The combination of the Neural-FGMRES algebraic preconditioner and the dual-scale topological formulation yields structural speedups exceeding $1,250\times$ (Figure~\ref{fig:serverless_economics}b). This enables hyper-local $\mathcal{O}(10)$ km global ensembles to be generated at sub-cent cost profiles, fundamentally democratizing access to verified high-resolution climate intelligence.
